% ── Dánský model flexicurity ─────────────────────────────────────────────────

Dánský model \acgen{KV} stojí na~třech vzájemně se podpírajících pilířích, které
ekonomická literatura souhrnně označuje jako~\emph{zlatý trojúhelník flexicurity}:
\begin{enumerate}[label=(\roman*)]
  \item vysoká flexibilita propouštění a~náboru bez zákonné minimální mzdy;
  \item štědrý systém podpor v~nezaměstnanosti čerpaný skrze odborové fondy (\textit{Ghent system});
  \item rozsáhlá aktivní politika zaměstnanosti zaměřená na~rekvalifikace.
\end{enumerate}
Trh práce není regulován zákonem o~mzdě ani o~délce pracovní doby --- veškeré
mzdové i~pracovní podmínky jsou upraveny výlučně \acs{KS}.
Smlouvy se uzavírají na~odvětvové úrovni mezi konfederacemi (\textit{LO},
nyní \textit{FH}, na~straně zaměstnanců a~\textit{DA} na~straně zaměstnavatelů)
a~jejich obsah je následně transponován do~podnikových dohod. Stát do~vyjednávání
nevstupuje --- jeho role se omezuje na~poskytování dávkového rámce
a~financování \acgen{APZ}.\par

Pokrytí \acs{KS} dosahuje v~Dánsku přibližně \SI{82}{\percent}
zaměstnanců (data 2018, viz obr.~\ref{fig:eu_pokryti_kv_mapa}), odborová organizovanost
$\approx \SI{67}{\percent}$ (obr.~\ref{fig:eu_hustota_mapa}) ji řadí
po~bok Švédska a~Finska na~špici \aclgen{EU}.
Vysoká odborová organizovanost není dána zákonnou povinností členství, ale
\textit{Ghent systémem}: dávky v~nezaměstnanosti administrují odborové fondy
(\textit{a-kasse}), členství v~nichž je~pro pracovníka racionální i~z~čistě
pojistného hlediska. Tento mechanismus udržuje vysokou
odborovou organizovanost i~v~období dlouhodobého poklesu odborářství v~ostatních
zemích~\acs{EU}. V~tomto modelu se~\acp{KS} uzavírají
na~tříletých cyklech a~mezisektorové vyjednávání zajišťuje koordinaci tak,
aby~mzdové nárůsty respektovaly konkurenceschopnost exportního sektoru.\par

Z~pohledu pracovního práva se~Dánsko vyznačuje absencí zákonné minimální mzdy
i~zákonné výpovědní doby pro~většinu profesí --- obě institucionální veličiny
jsou plně v~kompetenci sociálních partnerů. Propouštění z~organizačních důvodů
je~výrazně jednodušší než~v~kontinentální Evropě (krátké výpovědní lhůty,
nízké odstupné), což zvyšuje fluktuaci pracovní síly. Sociální dopad je~však
tlumen vysokou náhradovou mírou dávek (přibližně \SI{90}{\percent} mzdy
pro~nízkopříjmové pracovníky po~dobu až~dvou let) a~povinným zapojením
do~rekvalifikačního programu.
Celkový roční rozpočet \acgen{APZ} dosahuje hodnot blížících se~\SI{2}{\percent}
\acgen{HDP} --- nejvyšší úrovně v~\acs{OECD}.~\cite{OECD_GoodJobsForAll_2018}\par

Důsledkem této kombinace je~vysoká pracovní mobilita (přibližně 25~\%~zaměstnanců
mění zaměstnavatele ročně), nízká strukturální nezaměstnanost
a~zároveň jedny z~nejnižších hodnot Giniho koeficientu disponibilního příjmu
v~\acs{EU} ($\approx 0{,}27$). Mzdová komprese probíhá skrze koordinované
\acs{KV}, nikoli skrze redistribuční zdanění.\par

Přímý import dánského modelu naráží na~tři strukturální překážky.
\begin{enumerate}
  \item Za~prvé, ghentský systém vyžaduje legislativní rozhodnutí o~delegaci administrace
podpor v~nezaměstnanosti na~odbory --- politicky obtížné v~zemi s~odborovou organizovaností
pod \SI{12}{\percent}, kde~by~taková reforma neměla zaměstnaneckou
základnu.
  \item Za~druhé, absence zákonné minimální mzdy je~v~\acgen{ČR}
nepředstavitelná bez~předchozího dosažení vysokého pokrytí \acs{KS}
(viz~obecná podmínka směrnice~\acs{EU} 2022/2041, podle níž je~minimální mzda nezbytná v~zemích s~pokrytím pod \SI{80}{\percent}).~\cite{eu_directive_2022_2041}
  \item Za~třetí, fiskální náklady dánské \acgen{APZ} (\SI{2}{\percent} \acgen{HDP}) jsou
mnohonásobkem české úrovně (přibližně \SI{0,3}{\percent} \acgen{HDP}; viz
obr.~\ref{fig:eu_apz_vydaje_2004}).
\end{enumerate}
Hodnotu pro~\acacc{ČR} má~především princip
\emph{koordinovaného mezisektorového vyjednávání}, který lze přenést bez
nutnosti převzít celý fiskální rámec.~\cite{Denmark_LabourMarket}\par
